\documentclass[12pt]{article}
\usepackage[T1]{fontenc}
\usepackage[utf8]{inputenc}
\usepackage{lmodern}
\usepackage{textcomp}
\usepackage{enumitem}
\usepackage{geometry}
\geometry{margin=1in}
\begin{document}
\begin{center}
Answer Key - Geography - K-12 Level Assessment 4 \
\small (Answers are ordered exactly as in the question paper.)
\end{center}

\section*{Multiple Choice}
\begin{enumerate}[label=\arabic*.]
\item \textbf{Answer:} B
\\ \textit{Options:} A) dictatorship.; B) theocracy.; C) democracy.; D) autocracy.
\\ \textit{Note:} A theocracy is a system of government in which religious leaders control the government, and the state's legal system is based on religious law, effectively combining religion and governance.
\item \textbf{Answer:} A
\\ \textit{Options:} A) Zero population growth; B) High mortality rates; C) High birth rates; D) High sex ratios
\\ \textit{Note:} Zero population growth characterizes stages 4 and 5 in the demographic transition model because these stages are marked by low birth and death rates, leading to a stable population size where the number of births equals the number of deaths.
\item \textbf{Answer:} D
\\ \textit{Options:} A) Denmark; B) Finland; C) Luxembourg; D) Spain
\\ \textit{Note:} Spain is home to the Basque Country, where the Basque language, or Euskara, is co-official alongside Spanish, resulting in the highest concentration of Basque speakers in the world.
\item \textbf{Answer:} D
\\ \textit{Options:} A) regionalism.; B) diffusion.; C) iconography.; D) nationalism.
\\ \textit{Note:} Nationalism is the ideology that emphasizes loyalty and devotion to a nation-state, often leading individuals to identify with its culture, history, and goals.
\item \textbf{Answer:} A
\\ \textit{Options:} A) Hindu and Buddhist; B) Hindu and Muslim; C) Muslim and Jewish; D) Christian and Buddhist
\\ \textit{Note:} Hinduism and Buddhism both emphasize the belief in the cycle of rebirth and the impermanence of physical form, leading to the long-held tradition of cremation as a means of releasing the spirit from the body.
\end{enumerate}

\section*{True / False}
\begin{enumerate}[label=\arabic*.]
\setcounter{enumi}{5}
\item \textbf{Answer:} True
\\ \textit{Options:} A) True; B) False
\\ \textit{Note:} The statement is true because the majority of the Kurdish population resides in Turkey, where they constitute a substantial ethnic minority, particularly in the southeastern region of the country.
\item \textbf{Answer:} True
\\ \textit{Options:} A) True; B) False
\\ \textit{Note:} The statement is true because a world map encompasses a vast area, requiring a smaller scale to represent the entire globe accurately, while maps of smaller regions can use larger scales to provide more detailed information.
\item \textbf{Answer:} False
\\ \textit{Options:} A) True; B) False
\\ \textit{Note:} The statement is false because natural gas is actually the most abundant fossil fuel, and oil is primarily used for transportation and industrial purposes rather than just heating and cooking.
\item \textbf{Answer:} False
\\ \textit{Options:} A) True; B) False
\\ \textit{Note:} The statement is false because the Spanish were the first Europeans to establish settlements in the eastern United States, with St. Augustine, Florida, founded in 1565, predating French colonies in the region.
\item \textbf{Answer:} False
\\ \textit{Options:} A) True; B) False
\\ \textit{Note:} Heathenism traditionally refers to a belief system or practices associated with polytheism or paganism, rather than a movement away from religious beliefs and practices altogether.
\end{enumerate}

\section*{Long Form Response}
\begin{enumerate}[label=\arabic*.]
\setcounter{enumi}{10}
\item \textbf{Answer:} Universalizing religions, such as Christianity and Buddhism, are characterized by their appeal to a global audience, inclusive beliefs, and practices that transcend cultural boundaries. Christianity teaches the importance of love, forgiveness, and salvation, inviting followers from diverse backgrounds to share in its message. Similarly, Buddhism emphasizes practices like meditation and mindfulness, promoting peace and enlightenment for all individuals, regardless of their cultural origins. Both religions have spread widely across the globe, influencing societies and cultures, evident in the church's presence in various countries and Buddhist temples around the world. This global influence illustrates how universalizing religions can foster a sense of community and shared purpose among people from different backgrounds.
\\ \textit{Note:} The answer correctly identifies universalizing religions by highlighting their global appeal and inclusive beliefs, demonstrated through Christianity's message of love and salvation and Buddhism's practices of meditation and mindfulness, which resonate across diverse cultures.
\item \textbf{Answer:} Agriculture is the most widespread primary economic activity globally, serving as the foundation for many economies. It significantly impacts economies by providing jobs, generating income, and ensuring food security for growing populations. In societies, agriculture shapes cultural practices and social structures, as communities often rely on farming traditions and local produce. Environmentally, agriculture can both positively and negatively affect ecosystems; while it can promote biodiversity through sustainable practices, intensive farming can lead to deforestation, soil degradation, and loss of habitats. Therefore, understanding agriculture's role is essential for addressing economic, social, and environmental challenges worldwide.
\\ \textit{Note:} This answer is correct because it comprehensively addresses the multifaceted significance of agriculture, highlighting its essential role in economic development, social structure, and environmental impact, thereby reflecting its status as the most widespread primary economic activity globally.
\end{enumerate}

\vfill
\noindent\textit{End of Answer Key}
\end{document}